% ---- pg-420 （书页 408） ----

利用从两种解法估计出的载荷得到的因子得分配对图, 也是探测离群值的好工具. 若对某因子的载荷趋向一致, 离群值将在 $45^\circ$ 直线的邻域中出现, 但却远离原点和其余点聚集处. 从\hyperref[fig:9.6]{图 9.6} 的(b)图可以清楚看到, 276 个观测值中只有一个点与其他点不一致. 它有一个少见的大 $F_2$ 得分. 要是除去此点 $[39.1, 39.3, 75.7, 115, 73.4, 69.1]$, 并重作分析, 载荷不会有明显改变.

在数据集大时, 应将它分为两个(大致)相等的组, 并且在每个组上完成因子分析. 这些分析的结果可以彼此作比较, 还可以与全部数据集的分析作比较, 以检查解的稳定性. 如果结果相互一致, 对解的信任程度就增强了.

鸡的骨骼数据曾被分为 $n_1=137$ 个和 $n_2=139$ 个观察值的两组. 得到的样本相关矩阵为
\[
\mathbf{R}_1=\left[\begin{array}{rrrrrr}
1.000& & & & & \\
0.696&1.000& & & & \\
0.588&0.540&1.000& & & \\
0.639&0.575&0.901&1.000& & \\
0.694&0.606&0.844&0.835&1.000& \\
0.660&0.584&0.866&0.863&0.931&1.000
\end{array}\right]
\]
及
\[
\mathbf{R}_2=\left[\begin{array}{rrrrrr}
1.000& & & & & \\
0.366&1.000& & & & \\
0.572&0.352&1.000& & & \\
0.587&0.406&0.950&1.000& & \\
0.587&0.420&0.909&0.911&1.000& \\
0.598&0.386&0.894&0.927&0.940&1.000
\end{array}\right]
\]

旋转后的估计载荷、特殊方差及对 $m=3$ 因子模型主成分分解所解释的(标准化)样本总方差的比例, 在\hyperref[tab:9.11]{表 9.11} 中给出.

\begin{mvtable}
\phantomsection\label{tab:9.11}\mvtabcap{表 9.11}
\small
\begin{tabular}{lcccc|cccc}
\toprule
& \multicolumn{4}{c}{第一组 旋转因子载荷估计} & \multicolumn{4}{c}{第二组 旋转因子载荷估计}\\
& \multicolumn{4}{c}{($n_1=137$ 个观测值)} & \multicolumn{4}{c}{($n_2=139$ 个观测值)}\\
\cmidrule(r){2-5}\cmidrule(l){6-9}
变\quad 量 & $F_1^{*}$ & $F_2^{*}$ & $F_3^{*}$ & $\tilde{\psi}_i$ & $F_1^{*}$ & $F_2^{*}$ & $F_3^{*}$ & $\hat{\psi}_i$\\
\midrule
1. 颅骨长度 & 0.360 & 0.361 & 0.853 & 0.01 & 0.352 & 0.921 & 0.167 & 0.00\\
2. 颅骨幅宽 & 0.303 & 0.899 & 0.312 & 0.00 & 0.203 & 0.145 & 0.968 & 0.00\\
3. 股骨长 & 0.914 & 0.238 & 0.175 & 0.08 & 0.930 & 0.239 & 0.130 & 0.06\\
4. 胫骨长 & 0.877 & 0.270 & 0.242 & 0.10 & 0.925 & 0.248 & 0.187 & 0.05\\
5. 肱骨长 & 0.830 & 0.247 & 0.395 & 0.11 & 0.912 & 0.252 & 0.208 & 0.06\\
6. 尺骨长 & 0.871 & 0.231 & 0.332 & 0.08 & 0.914 & 0.272 & 0.168 & 0.06\\
\midrule
\makecell[l]{所解释的(标准化)\\样本总方差的\\累积比例}
 & 0.546 & 0.743 & 0.940 & & 0.593 & 0.780 & 0.962 & \\
\bottomrule
\end{tabular}
\end{mvtable}
